\chapter{CHAPTER 2: LITERATURE
REVIEW}\label{chapter-2-literature-review}

\section{2.1 Chapter Overview}\label{chapter-overview}

This chapter provides a comprehensive, critical synthesis of the
theoretical and empirical foundations governing air quality forecasting,
explicitly structured to address the latest advancements up to 2026.
Moving beyond chronological summaries, the literature is evaluated
thematically. Section 2.2 outlines the historical milestones of the
domain. Section 2.3 defines the core technical terminology and
frameworks. Section 2.4 critically analyzes topographical dispersion
dynamics, while Section 2.5 evaluates monitoring technologies. Section
2.6 traces the evolution of predictive algorithms. Section 2.7 provides
a direct comparative analysis of the three most relevant similar works
(W1, W2, W3), explicitly detailing their advantages and disadvantages to
establish the foundational rationale for the implementation detailed in
Chapter 5.

\section{2.2 History of the Domain (Key Milestones up to
2026)}\label{history-of-the-domain-key-milestones-up-to-2026}

The domain of air quality forecasting has undergone a radical
transformation over the past three decades, evolving from static
chemical dispersion models to dynamic, AI-driven predictive systems. *
\textbf{Pre-2010 (The Statistical Era):} Early forecasting relied
heavily on deterministic chemical transport models (CTMs) and linear
statistical time-series models like ARIMA. These models were
computationally expensive and required massive arrays of
government-grade sensors. * \textbf{2010--2018 (The Rise of Machine
Learning):} The proliferation of open data allowed researchers to apply
Support Vector Machines (SVM) and early Random Forest models to
meteorological data. However, accuracy was limited to 1-to-3 hour
horizons. * \textbf{2018--2022 (The Deep Learning Shift \& Low-Cost
Sensing):} The deployment of low-cost IoT sensor networks (like
PurpleAir) provided the high-frequency big data required to train
Recurrent Neural Networks (RNNs) and LSTMs. Simultaneously, Sentinel-5P
TROPOMI launched, democratizing global chemical satellite data. *
\textbf{2023--2026 (The Era of Explainability \& Multi-Modal Fusion):}
Modern literature now rejects ``black box'' neural networks in favor of
Explainable AI (SHAP). The current cutting-edge (as of 2026) focuses on
fusing ground telemetry with meteorological proxies using
computationally efficient tree-based algorithms (XGBoost), prioritizing
interpretability for immediate public health interventions.

\section{2.3 Technical Terminology and
Frameworks}\label{technical-terminology-and-frameworks}

To contextualize the comparative analyses in this chapter, the following
computational algorithms and frameworks are defined: * \textbf{SARIMAX
(Seasonal Autoregressive Integrated Moving Average with eXogenous
factors):} A traditional statistical baseline model that leverages
linear regression and historical memory to forecast time-series data,
accounting for seasonality and external covariates (e.g., temperature).
* \textbf{LSTM (Long Short-Term Memory):} A deep learning architecture
belonging to the RNN family. It utilizes internal ``gates'' (forget,
input, output) to selectively remember or discard temporal sequences,
theoretically making it ideal for long-term meteorological dependencies.
* \textbf{XGBoost (Extreme Gradient Boosting):} An advanced ensemble
machine learning framework based on decision trees. It sequentially
builds trees where each new tree corrects the residual errors of the
previous ones. It requires time-series data to be flattened into tabular
formats using engineered ``lag'' features. * \textbf{SHAP (SHapley
Additive exPlanations):} An Explainable AI (XAI) framework based on
cooperative game theory. It calculates ``Shapley values'' to isolate the
exact mathematical contribution (positive or negative) of a single
feature (e.g., wind speed) toward a specific algorithmic prediction,
effectively unlocking black-box models.

\section{2.4 Atmospheric Dynamics and Topographical
Influences}\label{atmospheric-dynamics-and-topographical-influences}

The dispersion of fine particulate matter (PM2.5) is inextricably linked
to topographical conditions (WHO, 2021). Recent studies demonstrate that
treating urban air quality homogeneously leads to predictive failure
(Dhammapala et al., 2022). Coastal environments like Colombo benefit
from the ``sea breeze effect,'' where land-sea thermal differentials
flush particulate matter inland (Kurukulasuriya et al., 2025). Peak
pollution here aligns with rush-hour traffic before rapid dispersion.
Conversely, valley basins like Kandy exhibit the ``valley trapping''
phenomenon. Thermal inversions act as a meteorological lid, trapping NO2
and PM2.5 within the basin for extended periods (Ambanwala, 2025).
\textbf{Link to Gap:} Existing predictive models for Sri Lanka treat
spatial data blindly. A direct empirical benchmarking of forecasting
accuracy between a coastal metropolis and a valley basin is missing from
the literature.

\section{2.5 Technological Approaches to Air Quality
Monitoring}\label{technological-approaches-to-air-quality-monitoring}

The literature reveals a bifurcation in monitoring technologies:
ground-based sensor networks (PurpleAir) and spaceborne remote sensing
(Sentinel-5P). While satellite data offers unparalleled spatial
coverage, it introduces severe temporal constraints (daily passes) and
is highly susceptible to cloud occlusion. In tropical, mountainous
regions like Kandy, cloud cover renders daily satellite vectors
fundamentally unreliable for high-frequency forecasting (Attanayake et
al., 2025). \textbf{Link to Gap:} The practical reality of cloud
occlusion in Sri Lanka makes satellite dependency a vulnerability. The
literature has not adequately proven whether high-frequency ground
telemetry augmented with localized weather APIs (Open-Meteo) can achieve
superior short-term accuracy compared to occlusion-prone satellite
models.

\section{2.6 Evolution of PM2.5 Forecasting
Methodologies}\label{evolution-of-pm2.5-forecasting-methodologies}

Traditional time-series models (SARIMAX) excel at capturing predictable
linear cycles but fail against sudden, non-linear pollution spikes
driven by erratic weather (Fernando et al., 2022). To resolve this, deep
learning (LSTM) became prominent. However, while LSTMs excel on massive
global datasets, they frequently overfit on sparse, volatile local
datasets. Consequently, Extreme Gradient Boosting (XGBoost) has emerged
as a formidable alternative. When provided with an engineered feature
vector of rolling statistical means and time-lags, tree-based models
capture non-linear relationships with greater computational efficiency
and lower susceptibility to overfitting than deep neural networks
(Kurukulasuriya et al., 2025). \textbf{Link to Gap:} The literature
presents conflicting evidence regarding the optimal algorithm for
multi-horizon environmental forecasting. A systematic benchmarking
pitting SARIMAX, LSTM, and XGBoost across continuously expanding
horizons (1 to 48 hours) is required.

\section{2.7 Similar Works Analysis (Comparative
Benchmarking)}\label{similar-works-analysis-comparative-benchmarking}

To establish a direct baseline for the implementation detailed in
Chapter 5, three seminal works within the localized and regional context
were analyzed. These represent the current state-of-the-art that the
SentinelAQ system aims to supersede.

\subsection{W1: Kurukulasuriya et al.~(2025) -- ``Uni-Modal Deep
Learning for Coastal
AQI''}\label{w1-kurukulasuriya-et-al.-2025-uni-modal-deep-learning-for-coastal-aqi}

\begin{itemize}
\tightlist
\item
  \textbf{Methodology:} Employed an LSTM architecture trained strictly
  on historical PM2.5 data from Colombo to predict 24-hour horizons.
\item
  \textbf{Advantages:} Achieved a low short-term RMSE by effectively
  mapping predictable diurnal traffic patterns in Colombo. It
  established that deep learning could be operationalized on low-cost
  sensor data.
\item
  \textbf{Disadvantages:} It is completely \emph{uni-modal} and
  \emph{spatially blind}. By excluding meteorological data (wind,
  humidity), the model systematically failed to predict transboundary
  smog events blowing in from the ocean. Furthermore, as an opaque
  neural network, it offered zero explainability for its predictions.
\end{itemize}

\subsection{W2: Dhammapala et al.~(2022) -- ``Static Aerosol Mapping via
Satellite''}\label{w2-dhammapala-et-al.-2022-static-aerosol-mapping-via-satellite}

\begin{itemize}
\tightlist
\item
  \textbf{Methodology:} Utilized MODIS Aerosol Optical Depth (AOD)
  satellite imagery to create an annual static mapping of PM2.5
  distribution across Sri Lanka.
\item
  \textbf{Advantages:} Provided unparalleled spatial coverage,
  identifying previously unknown pollution hotspots in
  rural/agricultural areas outside the purview of ground sensors.
\item
  \textbf{Disadvantages:} The approach is entirely \emph{static and
  retrospective}. It cannot forecast future events. Furthermore, the
  methodology explicitly cited massive data loss during the monsoon
  seasons due to severe cloud occlusion blocking the optical sensors,
  highlighting the fragility of satellite reliance in tropical climates.
\end{itemize}

\subsection{W3: Rowley \& Karakuş (2023) -- ``Global Multi-Modal
Ensemble
Forecasting''}\label{w3-rowley-karakuux15f-2023-global-multi-modal-ensemble-forecasting}

\begin{itemize}
\tightlist
\item
  \textbf{Methodology:} Fused ground sensor data with weather parameters
  using a Random Forest ensemble model to forecast PM2.5 in European
  cities.
\item
  \textbf{Advantages:} Proved that tree-based ensemble methods augmented
  with meteorological data outperform deep learning in multi-horizon
  forecasting, achieving high accuracy with lower computational
  overhead.
\item
  \textbf{Disadvantages:} The model was calibrated for temperate
  European climates with highly stable seasonal weather patterns. The
  specific hyperparameter tuning and feature selection lack
  applicability to the erratic, high-humidity, monsoon-driven dynamics
  of Sri Lankan microclimates. Additionally, it did not utilize SHAP for
  real-time feature attribution.
\end{itemize}

\textbf{Conclusion of Similar Works:} By synthesizing the advantages and
failures of W1, W2, and W3, it is evident that a successful architecture
for Sri Lanka must (a) avoid reliance on occluded satellite data (unlike
W2), (b) incorporate robust meteorological proxies to avoid spatial
blindness (unlike W1), and (c) deploy explainable, tree-based models
calibrated specifically for tropical monsoons (advancing W3). This
direct comparative analysis forms the evolutionary foundation for the
iterative implementation (v1 → v2 → v3) presented in Chapter 5.

\section{2.8 Summary and Identification of Research
Gaps}\label{summary-and-identification-of-research-gaps}

This critical review establishes that localized application of advanced
machine learning in Sri Lanka remains fragmented. Existing models (W1,
W2) operate as opaque black boxes or rely on occlusion-prone satellites.
By directly addressing the \emph{Multi-Horizon Predictive Gap}, the
\emph{Topographical Comparative Gap}, and the \emph{Interpretability
Gap}, this research synthesizes these disparate theoretical threads into
a singular, transparent, and highly accurate forecasting architecture.
The following chapters will detail the precise methodology,
requirements, and iterative development execution designed to overcome
the limitations of W1, W2, and W3.
